% ---
% title: Latent Consistency Models（LCM）
% date: 2025/12/22
% tag: Diffusion
% ---
扩散模型生成一张图要跑很多步去噪，推理成本高。Latent Consistency Models（LCMs）的做法是把扩散模型的概率流动力学蒸馏进一个学生模型，让它几步甚至一步出图。

\hrule

\section{扩散模型与概率流 ODE}

扩散模型的生成过程是从高斯噪声逐步还原到数据分布。写成连续形式，就是概率流常微分方程（Probability Flow ODE）：

\[
\frac{dx}{dt} = f_\theta(x, t)
\]

其中 \(x\) 是当前状态（带噪数据或 latent），\(t\) 是时间变量，\(f_\theta\) 是神经网络参数化的向量场。这个 ODE 的解就是从初始噪声 \(x_T\) 到目标数据 \(x_0\) 的确定性轨迹。

\hrule

\section{LCM 的建模目标}

LCM 的目标是学习一个映射函数：

\[
F_\theta(x_t, t) \approx x_0
\]

也就是从任意时间步 \(t\) 的状态 \(x_t\) 直接预测最终的无噪声结果 \(x_0\)。它不逐步积分 ODE，而是直接近似 ODE 的终点解。

\hrule

\section{一致性约束（Consistency Constraint）}

一致性约束要求同一个样本的不同噪声版本给出同一个预测：

\[
F_\theta(x_{t_1}, t_1) \approx F_\theta(x_{t_2}, t_2)
\]

其中 \(t_1, t_2\) 是不同时间步，对应同一数据样本的不同噪声版本。

\hrule

\section{蒸馏过程}

LCM 通过蒸馏预训练扩散模型（teacher）来训练学生模型（student）：

\begin{enumerate}
  \item 从 teacher 模型生成中间状态：
  
  \[
  x_t, x_{t'}
  \]
  \item 使用 teacher 模型计算目标输出：
  
  \[
  x_0^{\text{target}}
  \]
  \item 优化 student 模型：
  
  \[
  \mathcal{L}_{\text{distill}} = \|F_\theta(x_t, t) - x_0^{\text{target}}\|^2
  \]
  \item 同时加入一致性损失：
  
  \[
  \mathcal{L}_{\text{consistency}} = \|F_\theta(x_t, t) - F_\theta(x_{t'}, t')\|
  \]
\end{enumerate}

\hrule

\section{Latent Space 表示}

LCM 通常在 latent space 中进行建模，例如结合变分自编码器（VAE）：

\[
z = \mathcal{E}(x), \quad x = \mathcal{D}(z)
\]

其中 \(\mathcal{E}\) 是编码器，\(\mathcal{D}\) 是解码器。在 latent 空间建模能省掉大量计算。

\hrule

\section{推理过程}

LCM 推理过程可表示为：

\[
x_0 = F_\theta(x_T, T)
\]

在实践中，可采用少量离散时间步进行迭代更新：

\[
x_{t_{k-1}} = F_\theta(x_{t_k}, t_k)
\]

其中步数通常为 1–4 步。

\hrule

\section{方法对比}

\begin{center}
\begin{tabular}{llll}
\toprule
方法 & 建模目标 & 推理方式 & 推理步数 \\
\midrule
DDPM & 噪声预测 \(\epsilon\) & 逐步去噪 & 50–1000 \\
DDIM & 确定性轨迹 & 加速采样 & 10–50 \\
LCM & ODE 解近似 & 直接映射 & 1–4 \\
\bottomrule
\end{tabular}
\end{center}

\hrule

\section{理论解释}

LCM 把数值积分换成了函数逼近。原来的积分问题：

\[
x_0 = x_T + \int_T^0 f_\theta(x_t, t) dt
\]

变成了：

\[
F_\theta(x_t, t) \approx x_0
\]

也就是用神经网络直接近似 ODE 的积分结果。

\hrule

\section{优势与局限}

\subsection{优势}

推理快，能直接套在现有扩散模型上，也适合实时生成。

\subsection{局限}

代价是步数压到极少时画质会掉，效果依赖 teacher 模型的质量，复杂条件输入下稳定性有限。
